\chapter{The Yoneda Lemma}

\begin{overviewbox}
We come now to the most quoted single theorem of category theory:
the \term{Yoneda Lemma}. The statement is short. The consequences are
vast.

The starting point is the family of \term{hom-functors}. For each object
$A$ of $\mathcal{C}$, the functor $\mathrm{Hom}(A, -) : \mathcal{C} \to
\mathbf{Set}$ records ``how the rest of $\mathcal{C}$ maps out of $A$.''
This is a functor that captures $A$ \emph{from outside}: not what $A$ is
made of, but how it sits in the web of morphisms.

The Yoneda Lemma says, in one line:
\begin{equation*}
  \mathrm{Nat}(\mathrm{Hom}(A, -), F) \;\cong\; F(A).
\end{equation*}
Read: ``Natural transformations from $\mathrm{Hom}(A, -)$ into any functor
$F$ are in bijection with elements of $F(A)$.'' The bijection is given by
``evaluate the transformation at the identity $\mathrm{id}_A$.''

From this short statement flow many fundamental facts:
\begin{itemize}
  \item An object is completely determined by its hom-functor (up to iso).
        This is the \term{Yoneda embedding} of $\mathcal{C}$ into
        $[\mathcal{C}^{\mathrm{op}}, \mathbf{Set}]$.
  \item Universal properties can be cleanly recast as
        \term{representability}: a functor is representable iff it is
        isomorphic to some $\mathrm{Hom}(A, -)$.
  \item Many ``natural'' constructions in mathematics are, after Yoneda,
        just \emph{evaluation at the right object}.
\end{itemize}
We give the lemma, prove it, and then walk through its three signature
applications.
\end{overviewbox}


% ═══════════════════════════════════════════════════════════════════════════
\begin{nameorigin}
\term{Yoneda lemma} is named for \emph{Nobuo Yoneda} (1930--1996,
Japanese mathematician). The most-told story: Yoneda communicated the
result to Saunders Mac Lane on a Paris train platform during a 1954
mathematics conference; Mac Lane carried it back to the West. Yoneda
himself may never have published the lemma in its now-standard form —
he was primarily an algebraic topologist working on $\mathrm{Ext}$
groups, and the lemma was for him a tool, not a centerpiece. The lemma
has since become \emph{the} central organizing theorem of category
theory; appears in Vol I Ch.~9 (this), Vol II Ch.~15 (representability),
Vol III Ch.~28 (formal via ends), Vol IV Ch.~48 ($\infty$-categorical).

\term{Representable functor}: a functor isomorphic to some
$\mathrm{Hom}(A, -)$. The term ``representable'' is from
Eilenberg-Mac Lane; ``representing object'' is the object $A$ doing
the representing. The Yoneda philosophy: an object is determined by
\emph{what maps into it} or \emph{out of it}.
\end{nameorigin}

\section{Hom-Functors and Representability}
\label{sec:hom-functors}
% ═══════════════════════════════════════════════════════════════════════════

\subsection{Informally: An Object Seen Through Its Outgoing Maps}

Fix an object $A$ in $\mathcal{C}$. For each object $B$, the hom-set
$\mathrm{Hom}(A, B)$ tells us ``how many ways there are to map $A$ into
$B$, and what they look like.'' As $B$ varies, these hom-sets fit
together into a functor $\mathrm{Hom}(A, -) : \mathcal{C} \to
\mathbf{Set}$, with action on a morphism $g : B \to C$ given by
postcomposition: $g_* (f) = g \circ f$.

The functor $\mathrm{Hom}(A, -)$ \emph{is} $A$, seen from outside.
Everything about how $A$ relates to other objects is recorded there.

\subsection{Expanding: The Two Hom-Functors}

\begin{example}{Covariant Hom-Functor}
$\mathrm{Hom}(A, -) : \mathcal{C} \to \mathbf{Set}$ sends:
\begin{itemize}
  \item Object $B$ to the set $\mathrm{Hom}(A, B)$.
  \item Morphism $g : B \to C$ to the function
        $g_* = (g \circ -) : \mathrm{Hom}(A, B) \to \mathrm{Hom}(A, C)$.
\end{itemize}
Functoriality is immediate: $(g \circ f)_* = g_* \circ f_*$ and
$(\mathrm{id})_* = \mathrm{id}$, both by associativity and identity laws.
\end{example}

\begin{example}{Contravariant Hom-Functor}
$\mathrm{Hom}(-, A) : \mathcal{C}^{\mathrm{op}} \to \mathbf{Set}$ sends:
\begin{itemize}
  \item Object $B$ to the set $\mathrm{Hom}(B, A)$.
  \item Morphism $g : B \to C$ (in $\mathcal{C}$, reversed in
        $\mathcal{C}^{\mathrm{op}}$) to the function $g^* = (- \circ g)
        : \mathrm{Hom}(C, A) \to \mathrm{Hom}(B, A)$, given by
        \emph{pre}composition.
\end{itemize}
\end{example}

\begin{dialog}
Two hom-functors per object, one covariant and one contravariant. The
covariant one is ``maps out of $A$''; the contravariant is ``maps into
$A$.''
\end{dialog}

\subsection{Formalizing: Representable Functors}

A functor $F : \mathcal{C} \to \mathbf{Set}$ is \term{representable} if
there exists an object $A$ of $\mathcal{C}$ and a natural isomorphism
\begin{equation*}
  F \;\cong\; \mathrm{Hom}(A, -).
\end{equation*}
The object $A$ is called a \term{representing object} for $F$ (it is
unique up to isomorphism when it exists).

Dually, $F : \mathcal{C}^{\mathrm{op}} \to \mathbf{Set}$ is representable
if $F \cong \mathrm{Hom}(-, A)$ for some $A$.

\begin{howtosay}
\begin{itemize}
  \item ``representable'' --- the functor is ``representable by $A$.''
  \item $\mathrm{Hom}(A, -)$ --- ``the covariant hom out of $A$,''
        ``$A$'s functor of points,'' or simply ``the representable at $A$.''
  \item $\mathrm{Hom}(-, A)$ --- ``the contravariant hom into $A$,'' or
        ``the representable presheaf at $A$.''
\end{itemize}
\end{howtosay}

\subsection{Reorganizing: Universal Properties as Representability}

Most universal properties can be repackaged as ``the functor of $X$ with
such-and-such property is representable.'' For example: a product
$A \times B$ represents the functor $X \mapsto \mathrm{Hom}(X, A) \times
\mathrm{Hom}(X, B)$. A pullback represents the functor of pairs of
morphisms agreeing under composition. This packaging is uniform, and it
is the door to the Yoneda Lemma.

\subsection{Simplifying: The Hom-Functor Recipe}

To check whether $F : \mathcal{C} \to \mathbf{Set}$ is representable, look
for an $A$ such that $F(B) \cong \mathrm{Hom}(A, B)$ in a way natural in $B$.
If such an $A$ exists, $F$ is representable; the universal element of $F$
is the image of $\mathrm{id}_A$ in $F(A)$ under the corresponding isomorphism.

\subsection{Formally: Naturality of the Hom-Functor}

\formalnote{$\mathrm{Hom}$ as a bifunctor}
The two-sided $\mathrm{Hom}(-, -) : \mathcal{C}^{\mathrm{op}} \times
\mathcal{C} \to \mathbf{Set}$ is a bifunctor: contravariant in the first
slot, covariant in the second. The covariant and contravariant
hom-functors are obtained by fixing the appropriate slot.


% ═══════════════════════════════════════════════════════════════════════════
\section{The Yoneda Lemma}
\label{sec:yoneda-statement}
% ═══════════════════════════════════════════════════════════════════════════

\subsection{Informally: A Natural Transformation Is Just an Element}

Suppose we are given a functor $F : \mathcal{C} \to \mathbf{Set}$ and an
object $A$ of $\mathcal{C}$. How many natural transformations $\alpha :
\mathrm{Hom}(A, -) \Rightarrow F$ are there?

Pick any natural transformation $\alpha$. Its component at $A$ is a
function $\alpha_A : \mathrm{Hom}(A, A) \to F(A)$. Apply it to the special
element $\mathrm{id}_A \in \mathrm{Hom}(A, A)$:
\begin{equation*}
  x = \alpha_A(\mathrm{id}_A) \in F(A).
\end{equation*}
The Yoneda Lemma says: this single element $x$ determines $\alpha$
\emph{completely}, and conversely every $x \in F(A)$ arises this way from
a unique $\alpha$.

\subsection{Expanding: The Bijection}

Given $x \in F(A)$, define $\alpha^x : \mathrm{Hom}(A, -) \Rightarrow F$
by:
\begin{equation*}
  \alpha^x_B(f : A \to B) := F(f)(x) \in F(B).
\end{equation*}
That is, transport $x$ from $F(A)$ to $F(B)$ along $F(f)$.

The two assignments
\begin{align*}
  \alpha \;\mapsto\; \alpha_A(\mathrm{id}_A) \quad\text{and}\quad
  x \;\mapsto\; \alpha^x
\end{align*}
are inverse to each other. They establish a bijection between natural
transformations $\mathrm{Hom}(A, -) \Rightarrow F$ and elements of $F(A)$.

\subsection{Formalizing: The Statement}

\textbf{The Yoneda Lemma.}\quad For any functor $F : \mathcal{C} \to
\mathbf{Set}$ and any object $A$ of $\mathcal{C}$, there is a bijection
\begin{equation*}
  \boxed{\;
    \mathrm{Nat}(\mathrm{Hom}(A, -),\; F)
    \;\xrightarrow{\;\cong\;}\;
    F(A),
    \quad
    \alpha \;\mapsto\; \alpha_A(\mathrm{id}_A).
  \;}
\end{equation*}
This bijection is natural in both $A$ (contravariantly) and $F$ (covariantly).

\textbf{Proof.}\quad The map ``$\alpha \mapsto \alpha_A(\mathrm{id}_A)$''
is well-defined. We exhibit an inverse: given $x \in F(A)$, define
$\alpha^x_B(f) := F(f)(x)$. Naturality of $\alpha^x$ comes for free from
functoriality of $F$:
\begin{align*}
  \alpha^x_C(g \circ f)
    &= F(g \circ f)(x)
       && \text{(definition)} \\
    &= F(g)(F(f)(x))
       && \text{(functoriality)} \\
    &= F(g)(\alpha^x_B(f))
       && \text{(definition)}.
\end{align*}
This is the naturality square at the morphism $g : B \to C$ applied to
$f$.

To see ``$\alpha \mapsto \alpha_A(\mathrm{id}_A) \mapsto \alpha^{\alpha_A(\mathrm{id}_A)}$''
is the identity, fix any $f : A \to B$. We need
$\alpha^{\alpha_A(\mathrm{id}_A)}_B(f) = \alpha_B(f)$:
\begin{align*}
  \alpha^{\alpha_A(\mathrm{id}_A)}_B(f)
    &= F(f)(\alpha_A(\mathrm{id}_A))
       && \text{(definition of $\alpha^{(-)}$)} \\
    &= \alpha_B(f \circ \mathrm{id}_A)
       && \text{(naturality of $\alpha$ at $f$)} \\
    &= \alpha_B(f).
\end{align*}
The reverse composite ``$x \mapsto \alpha^x \mapsto \alpha^x_A(\mathrm{id}_A)$''
gives $F(\mathrm{id}_A)(x) = x$. So we have inverse bijections. \qed

\subsection{Reorganizing: The Contravariant Version}

Dually, for a contravariant functor $F : \mathcal{C}^{\mathrm{op}} \to
\mathbf{Set}$:
\begin{equation*}
  \mathrm{Nat}(\mathrm{Hom}(-, A),\; F) \;\cong\; F(A).
\end{equation*}
The bijection again sends $\alpha$ to $\alpha_A(\mathrm{id}_A)$.

\subsection{Simplifying: The One-Line Slogan}

\begin{equation*}
  \boxed{\;\;
    \mathrm{Nat}(\mathrm{Hom}(A,-), F) \;\cong\; F(A)
    \;\;\text{(natural in $A$, $F$).}
  \;\;}
\end{equation*}

\subsection{Formally: Naturality of the Yoneda Isomorphism}

\formalnote{Naturality in $F$ and $A$}
The bijection of the Yoneda Lemma is natural in two variables. For a
morphism $f : A' \to A$ in $\mathcal{C}$, we get a natural transformation
$f^* : \mathrm{Hom}(A, -) \Rightarrow \mathrm{Hom}(A', -)$ given by
precomposition, and the Yoneda bijection respects this. For a natural
transformation $\theta : F \Rightarrow G$, we similarly get compatible
maps. These coherences are what allow Yoneda to upgrade to the Yoneda
\emph{embedding}.


% ═══════════════════════════════════════════════════════════════════════════
\section{The Yoneda Embedding}
\label{sec:yoneda-embedding}
% ═══════════════════════════════════════════════════════════════════════════

\subsection{Informally: A Category Inside a Bigger Category}

Apply Yoneda with $F = \mathrm{Hom}(B, -)$ for a second object $B$:
\begin{equation*}
  \mathrm{Nat}(\mathrm{Hom}(A, -),\; \mathrm{Hom}(B, -)) \;\cong\; \mathrm{Hom}(B, A).
\end{equation*}
So natural transformations between the hom-functors of $A$ and $B$ are
exactly the morphisms $B \to A$ in $\mathcal{C}$. The hom-functors faithfully
encode $\mathcal{C}$.

This gives us a functor: take each object $A$ to its (contravariant)
hom-functor $\mathrm{Hom}(-, A)$. The Yoneda Lemma says this functor is
fully faithful. We have embedded $\mathcal{C}$ inside the much larger
category $[\mathcal{C}^{\mathrm{op}}, \mathbf{Set}]$ of presheaves on
$\mathcal{C}$.

\subsection{Expanding: The Embedding in Detail}

Define
\begin{equation*}
  \mathit{Y} : \mathcal{C} \to [\mathcal{C}^{\mathrm{op}}, \mathbf{Set}]
\end{equation*}
by $\mathit{Y}(A) = \mathrm{Hom}(-, A)$ on objects, and
$\mathit{Y}(g : A \to A') = (- \circ g)_*$ on morphisms (postcomposition
with $g$ in each component).

By the Yoneda Lemma applied to $F = \mathrm{Hom}(-, A')$:
\begin{equation*}
  \mathrm{Hom}_{[\mathcal{C}^{\mathrm{op}}, \mathbf{Set}]}(\mathit{Y}(A), \mathit{Y}(A'))
  \;\cong\; \mathrm{Hom}(-, A')(A) = \mathrm{Hom}(A, A').
\end{equation*}
So $\mathit{Y}$ is fully faithful: it gives a bijection between morphisms
in $\mathcal{C}$ and natural transformations between the representable
presheaves. Hence $\mathcal{C}$ sits inside
$[\mathcal{C}^{\mathrm{op}}, \mathbf{Set}]$ as a full subcategory, via the
Yoneda embedding.

\subsection{Formalizing: The Embedding Theorem}

\textbf{Theorem (Yoneda Embedding).}\quad The functor
\begin{equation*}
  \mathit{Y} : \mathcal{C} \to [\mathcal{C}^{\mathrm{op}}, \mathbf{Set}],
  \quad
  A \mapsto \mathrm{Hom}(-, A)
\end{equation*}
is fully faithful. It identifies $\mathcal{C}$ with the full subcategory
of representable presheaves.

\subsection{Reorganizing: Three Consequences}

\textbf{1. Objects are determined by their hom-functors.}\quad
If $\mathrm{Hom}(-, A) \cong \mathrm{Hom}(-, B)$ as presheaves, then
$A \cong B$. (A fully faithful functor reflects isomorphism.)

\textbf{2. Universal properties name objects up to canonical iso.}\quad
A universal property describes an object as the representing object of a
functor. By Yoneda, this object is unique up to canonical iso.

\textbf{3. ``$\mathbf{Set}$ is the universe of presheaves.''}\quad
$[\mathcal{C}^{\mathrm{op}}, \mathbf{Set}]$ is the ``free completion'' of
$\mathcal{C}$ under colimits, in a precise sense. Every presheaf is a
colimit of representables.

\subsection{Simplifying: The Yoneda Slogan}

\begin{equation*}
  \boxed{\;\;
    \mathcal{C} \hookrightarrow [\mathcal{C}^{\mathrm{op}}, \mathbf{Set}],
    \;\;
    A \mapsto \mathrm{Hom}(-, A).
  \;\;}
\end{equation*}

\subsection{Formally: Presheaves and Probing}

\formalnote{Presheaves}
A \term{presheaf on $\mathcal{C}$} is a functor $\mathcal{C}^{\mathrm{op}}
\to \mathbf{Set}$. Yoneda says $\mathcal{C}$ embeds in its category of
presheaves. The full subcategory of representable presheaves is
$\mathit{Y}(\mathcal{C})$, equivalent to $\mathcal{C}$ itself.

\formalnote{Probing with representables}
For any presheaf $F$, the Yoneda Lemma says $F(A) \cong
\mathrm{Nat}(\mathit{Y} A, F)$. So evaluating $F$ at $A$ is the same as
``probing'' $F$ with the representable at $A$. This is the perspective
that motivates the slogan ``points of a sheaf are maps from a
representable.''


% ═══════════════════════════════════════════════════════════════════════════
\section{Applications of Yoneda}
\label{sec:yoneda-apps}
% ═══════════════════════════════════════════════════════════════════════════

\subsection{Informally: Three Patterns, One Theorem}

We have three signature uses of Yoneda. Each is a translation: ``a
property of an object'' becomes ``a property of the corresponding
representable.''

\subsection{Expanding: Three Applications}

\begin{example}{Application 1: Adjunctions Restated}
$L \dashv R$ iff $\mathrm{Hom}(L A, B) \cong \mathrm{Hom}(A, R B)$
naturally. By Yoneda, this is the same as a natural isomorphism
$\mathrm{Hom}(L A, -) \cong \mathrm{Hom}(A, R -)$, or alternatively
$\mathrm{Hom}(-, R B) \cong \mathrm{Hom}(L -, B)$. So adjunction is a
relation between hom-functors, exposed by Yoneda.
\end{example}

\begin{example}{Application 2: Equality of Functors}
To show two functors $F, G : \mathcal{C} \to \mathcal{D}$ are isomorphic,
it suffices to show $\mathrm{Hom}(-, F(A)) \cong \mathrm{Hom}(-, G(A))$
naturally in $A$. By Yoneda, this is equivalent to $F(A) \cong G(A)$
canonically, hence to $F \cong G$ as functors. Many ``$F = G$'' arguments
in practice are Yoneda arguments in disguise.
\end{example}

\begin{example}{Application 3: Representability Lifts}
If $F$ is representable, then $F$ preserves all (small) limits that exist
in $\mathcal{C}$. (This is a special case of the fact from Chapter~6 that
$\mathrm{Hom}(A, -)$ preserves limits.) So representability is a strong
constraint: it forces a functor to be limit-preserving.
\end{example}

\subsection{Formalizing: Density of Representables}

\textbf{Theorem (Density).}\quad Every presheaf $F : \mathcal{C}^{\mathrm{op}}
\to \mathbf{Set}$ is canonically a colimit of representables. More
specifically, $F$ is the colimit of the diagram
\begin{equation*}
  \mathrm{el}(F) \to \mathcal{C} \xrightarrow{\mathit{Y}}
  [\mathcal{C}^{\mathrm{op}}, \mathbf{Set}],
\end{equation*}
where $\mathrm{el}(F)$ is the category of elements of $F$ (pairs $(A, x)$
with $x \in F(A)$).

\subsection{Reorganizing: The Yoneda Toolkit}

To use Yoneda effectively:
\begin{enumerate}
  \item Translate the problem about objects of $\mathcal{C}$ into a
        problem about representable presheaves.
  \item Use the embedding's full-faithfulness or the Yoneda bijection.
  \item Translate back if needed.
\end{enumerate}

This is the move ``check on representables.'' It is sometimes called
\emph{the Yoneda manoeuvre} or simply \emph{by Yoneda}.

\subsection{Simplifying: The Yoneda Maxim}

\begin{quote}
\textit{To know an object, know its hom-functor. To know a functor, know
how it acts on representables. To compare two things, compare them on
representables.}
\end{quote}

\subsection{Formally: The Co-Yoneda Lemma}

\formalnote{Co-Yoneda}
The dual statement: for a presheaf $F$, the canonical map from the colimit
$\int^A F(A) \times \mathrm{Hom}(-, A)$ to $F$ is an isomorphism. This is
the formal statement of ``every presheaf is a colimit of representables.''
The notation $\int^A$ here is a coend, a kind of co-limit indexed over a
``twisted'' diagram. Coends are beyond the scope of this chapter.


% ═══════════════════════════════════════════════════════════════════════════
\section*{Exercises}
\addcontentsline{toc}{section}{Exercises}
% ═══════════════════════════════════════════════════════════════════════════

\begin{exercisesbox}
\begin{qa}
\qaitem{What is $\mathrm{Hom}(A, -)$?}{The functor $\mathcal{C} \to \mathbf{Set}$ sending $B$ to the set of morphisms $A \to B$.}
\qaitem{Is $\mathrm{Hom}(A, -)$ covariant or contravariant?}{Covariant.}
\qaitem{What is $\mathrm{Hom}(-, A)$?}{The functor $\mathcal{C}^{\mathrm{op}} \to \mathbf{Set}$ sending $B$ to morphisms $B \to A$.}
\qaitem{How does $\mathrm{Hom}(A, -)$ act on $g : B \to C$?}{Postcomposition: $f \mapsto g \circ f$.}
\qaitem{How does $\mathrm{Hom}(-, A)$ act on $g : B \to C$?}{Precomposition: $f \mapsto f \circ g$.}
\qaitem{What does it mean for $F$ to be representable?}{$F \cong \mathrm{Hom}(A, -)$ for some $A$.}
\qaitem{Is the representing object unique?}{Yes, up to isomorphism.}
\qaitem{State the Yoneda Lemma.}{$\mathrm{Nat}(\mathrm{Hom}(A, -), F) \cong F(A)$.}
\qaitem{What is the explicit bijection?}{$\alpha \mapsto \alpha_A(\mathrm{id}_A)$.}
\qaitem{What is the inverse map?}{$x \mapsto \alpha^x$ where $\alpha^x_B(f) = F(f)(x)$.}
\qaitem{Why is the inverse $\alpha^x$ natural?}{Because $F$ is a functor: $F(g \circ f) = F(g) F(f)$.}
\qaitem{What does $\alpha_A(\mathrm{id}_A)$ represent?}{The ``universal element'' that determines $\alpha$.}
\qaitem{Yoneda for contravariant functors: state it.}{$\mathrm{Nat}(\mathrm{Hom}(-, A), F) \cong F(A)$ for $F : \mathcal{C}^{\mathrm{op}} \to \mathbf{Set}$.}
\qaitem{What is the Yoneda embedding?}{The functor $A \mapsto \mathrm{Hom}(-, A)$ from $\mathcal{C}$ into $[\mathcal{C}^{\mathrm{op}}, \mathbf{Set}]$.}
\qaitem{Is the Yoneda embedding faithful?}{Yes.}
\qaitem{Is it full?}{Yes (this is what the Yoneda Lemma supplies).}
\qaitem{So the Yoneda embedding is what?}{Fully faithful.}
\qaitem{What is a presheaf on $\mathcal{C}$?}{A functor $\mathcal{C}^{\mathrm{op}} \to \mathbf{Set}$.}
\qaitem{What is a representable presheaf?}{One of the form $\mathrm{Hom}(-, A)$.}
\qaitem{Consequence: when are two objects isomorphic?}{Iff their representables are isomorphic as presheaves.}
\qaitem{Hom-functors preserve what?}{All small limits.}
\qaitem{Therefore representable functors preserve what?}{All small limits.}
\qaitem{What does the density theorem say?}{Every presheaf is canonically a colimit of representables.}
\qaitem{What is the slogan for using Yoneda?}{``Check on representables.''}
\qaitem{Yoneda restatement of adjunction $L \dashv R$?}{$\mathrm{Hom}(L A, -) \cong \mathrm{Hom}(A, R -)$ as natural transformations.}
\qaitem{Yoneda restatement of $A \cong B$?}{$\mathrm{Hom}(A, -) \cong \mathrm{Hom}(B, -)$ naturally.}
\qaitem{What is the ``category of elements'' of a presheaf $F$?}{The category with objects $(A, x)$, $x \in F(A)$, and morphisms $f : A \to A'$ with $F(f)(x') = x$.}
\qaitem{Why is the Yoneda Lemma considered ``deep''?}{It links naturality, representability, and elements in a single short statement.}
\qaitem{What was the most surprising consequence to you?}{Often: that universal properties identify objects up to unique iso \emph{via Yoneda}.}
\qaitem{Does Yoneda hold for general $\mathcal{V}$-enriched categories?}{Yes; it generalizes essentially verbatim, with $\mathbf{Set}$ replaced by $\mathcal{V}$.}
\end{qa}
\end{exercisesbox}


% ═══════════════════════════════════════════════════════════════════════════
\section*{Chapter Summary}
\addcontentsline{toc}{section}{Chapter Summary}
% ═══════════════════════════════════════════════════════════════════════════

\begin{summarybox}
The \textbf{hom-functors} $\mathrm{Hom}(A, -)$ and $\mathrm{Hom}(-, A)$
record how an object $A$ relates to all others. A functor is
\textbf{representable} if it is naturally isomorphic to such a hom-functor.

The \textbf{Yoneda Lemma}: natural transformations from
$\mathrm{Hom}(A, -)$ to any functor $F$ are in bijection with elements
of $F(A)$. Explicitly: $\alpha \mapsto \alpha_A(\mathrm{id}_A)$.

The \textbf{Yoneda Embedding} $A \mapsto \mathrm{Hom}(-, A)$ is fully
faithful, embedding $\mathcal{C}$ into the presheaf category
$[\mathcal{C}^{\mathrm{op}}, \mathbf{Set}]$.

Consequence: objects are determined by their hom-functors (up to iso),
representable functors preserve limits, and every presheaf is a colimit
of representables (\textbf{density}).

The slogan: \textbf{check on representables}.
\end{summarybox}


% ═══════════════════════════════════════════════════════════════════════════
\section*{Formal Restatement}
\addcontentsline{toc}{section}{Formal Restatement}
% ═══════════════════════════════════════════════════════════════════════════

\begin{formalrestatement}
\textbf{Hom-functor.}\quad
$\mathrm{Hom}(A, -) : \mathcal{C} \to \mathbf{Set}$ sends $B$ to
$\mathrm{Hom}(A, B)$ and $g : B \to C$ to $g_* = (g \circ -)$.

\medskip
\textbf{Representability.}\quad
$F : \mathcal{C} \to \mathbf{Set}$ is representable iff
$F \cong \mathrm{Hom}(A, -)$ for some $A$.

\medskip
\textbf{Yoneda Lemma.}\quad
For $F : \mathcal{C} \to \mathbf{Set}$ and $A \in \mathcal{C}$:
\begin{equation*}
  \mathrm{Nat}(\mathrm{Hom}(A, -),\; F) \;\cong\; F(A),
  \quad \alpha \;\mapsto\; \alpha_A(\mathrm{id}_A).
\end{equation*}
The bijection is natural in $A$ and $F$.

\medskip
\textbf{Yoneda Embedding.}\quad
$\mathit{Y} : \mathcal{C} \to [\mathcal{C}^{\mathrm{op}}, \mathbf{Set}]$,
$A \mapsto \mathrm{Hom}(-, A)$. Fully faithful.

\medskip
\textbf{Density.}\quad
Every presheaf $F$ is a colimit of representables, indexed by
$\mathrm{el}(F)$.

\medskip
\noindent\textit{Chapter~10 uses representability and Yoneda to develop
the Adjoint Functor Theorems --- general existence theorems for left and
right adjoints.}
\end{formalrestatement}
